% Options for packages loaded elsewhere
\PassOptionsToPackage{unicode}{hyperref}
\PassOptionsToPackage{hyphens}{url}
\documentclass[
]{article}
\usepackage{xcolor}
\usepackage[margin=1in]{geometry}
\usepackage{amsmath,amssymb}
\setcounter{secnumdepth}{-\maxdimen} % remove section numbering
\usepackage{iftex}
\ifPDFTeX
  \usepackage[T1]{fontenc}
  \usepackage[utf8]{inputenc}
  \usepackage{textcomp} % provide euro and other symbols
\else % if luatex or xetex
  \usepackage{unicode-math} % this also loads fontspec
  \defaultfontfeatures{Scale=MatchLowercase}
  \defaultfontfeatures[\rmfamily]{Ligatures=TeX,Scale=1}
\fi
\usepackage{lmodern}
\ifPDFTeX\else
  % xetex/luatex font selection
\fi
% Use upquote if available, for straight quotes in verbatim environments
\IfFileExists{upquote.sty}{\usepackage{upquote}}{}
\IfFileExists{microtype.sty}{% use microtype if available
  \usepackage[]{microtype}
  \UseMicrotypeSet[protrusion]{basicmath} % disable protrusion for tt fonts
}{}
\makeatletter
\@ifundefined{KOMAClassName}{% if non-KOMA class
  \IfFileExists{parskip.sty}{%
    \usepackage{parskip}
  }{% else
    \setlength{\parindent}{0pt}
    \setlength{\parskip}{6pt plus 2pt minus 1pt}}
}{% if KOMA class
  \KOMAoptions{parskip=half}}
\makeatother
\usepackage{color}
\usepackage{fancyvrb}
\newcommand{\VerbBar}{|}
\newcommand{\VERB}{\Verb[commandchars=\\\{\}]}
\DefineVerbatimEnvironment{Highlighting}{Verbatim}{commandchars=\\\{\}}
% Add ',fontsize=\small' for more characters per line
\usepackage{framed}
\definecolor{shadecolor}{RGB}{248,248,248}
\newenvironment{Shaded}{\begin{snugshade}}{\end{snugshade}}
\newcommand{\AlertTok}[1]{\textcolor[rgb]{0.94,0.16,0.16}{#1}}
\newcommand{\AnnotationTok}[1]{\textcolor[rgb]{0.56,0.35,0.01}{\textbf{\textit{#1}}}}
\newcommand{\AttributeTok}[1]{\textcolor[rgb]{0.13,0.29,0.53}{#1}}
\newcommand{\BaseNTok}[1]{\textcolor[rgb]{0.00,0.00,0.81}{#1}}
\newcommand{\BuiltInTok}[1]{#1}
\newcommand{\CharTok}[1]{\textcolor[rgb]{0.31,0.60,0.02}{#1}}
\newcommand{\CommentTok}[1]{\textcolor[rgb]{0.56,0.35,0.01}{\textit{#1}}}
\newcommand{\CommentVarTok}[1]{\textcolor[rgb]{0.56,0.35,0.01}{\textbf{\textit{#1}}}}
\newcommand{\ConstantTok}[1]{\textcolor[rgb]{0.56,0.35,0.01}{#1}}
\newcommand{\ControlFlowTok}[1]{\textcolor[rgb]{0.13,0.29,0.53}{\textbf{#1}}}
\newcommand{\DataTypeTok}[1]{\textcolor[rgb]{0.13,0.29,0.53}{#1}}
\newcommand{\DecValTok}[1]{\textcolor[rgb]{0.00,0.00,0.81}{#1}}
\newcommand{\DocumentationTok}[1]{\textcolor[rgb]{0.56,0.35,0.01}{\textbf{\textit{#1}}}}
\newcommand{\ErrorTok}[1]{\textcolor[rgb]{0.64,0.00,0.00}{\textbf{#1}}}
\newcommand{\ExtensionTok}[1]{#1}
\newcommand{\FloatTok}[1]{\textcolor[rgb]{0.00,0.00,0.81}{#1}}
\newcommand{\FunctionTok}[1]{\textcolor[rgb]{0.13,0.29,0.53}{\textbf{#1}}}
\newcommand{\ImportTok}[1]{#1}
\newcommand{\InformationTok}[1]{\textcolor[rgb]{0.56,0.35,0.01}{\textbf{\textit{#1}}}}
\newcommand{\KeywordTok}[1]{\textcolor[rgb]{0.13,0.29,0.53}{\textbf{#1}}}
\newcommand{\NormalTok}[1]{#1}
\newcommand{\OperatorTok}[1]{\textcolor[rgb]{0.81,0.36,0.00}{\textbf{#1}}}
\newcommand{\OtherTok}[1]{\textcolor[rgb]{0.56,0.35,0.01}{#1}}
\newcommand{\PreprocessorTok}[1]{\textcolor[rgb]{0.56,0.35,0.01}{\textit{#1}}}
\newcommand{\RegionMarkerTok}[1]{#1}
\newcommand{\SpecialCharTok}[1]{\textcolor[rgb]{0.81,0.36,0.00}{\textbf{#1}}}
\newcommand{\SpecialStringTok}[1]{\textcolor[rgb]{0.31,0.60,0.02}{#1}}
\newcommand{\StringTok}[1]{\textcolor[rgb]{0.31,0.60,0.02}{#1}}
\newcommand{\VariableTok}[1]{\textcolor[rgb]{0.00,0.00,0.00}{#1}}
\newcommand{\VerbatimStringTok}[1]{\textcolor[rgb]{0.31,0.60,0.02}{#1}}
\newcommand{\WarningTok}[1]{\textcolor[rgb]{0.56,0.35,0.01}{\textbf{\textit{#1}}}}
\usepackage{graphicx}
\makeatletter
\newsavebox\pandoc@box
\newcommand*\pandocbounded[1]{% scales image to fit in text height/width
  \sbox\pandoc@box{#1}%
  \Gscale@div\@tempa{\textheight}{\dimexpr\ht\pandoc@box+\dp\pandoc@box\relax}%
  \Gscale@div\@tempb{\linewidth}{\wd\pandoc@box}%
  \ifdim\@tempb\p@<\@tempa\p@\let\@tempa\@tempb\fi% select the smaller of both
  \ifdim\@tempa\p@<\p@\scalebox{\@tempa}{\usebox\pandoc@box}%
  \else\usebox{\pandoc@box}%
  \fi%
}
% Set default figure placement to htbp
\def\fps@figure{htbp}
\makeatother
\setlength{\emergencystretch}{3em} % prevent overfull lines
\providecommand{\tightlist}{%
  \setlength{\itemsep}{0pt}\setlength{\parskip}{0pt}}
\usepackage{bookmark}
\IfFileExists{xurl.sty}{\usepackage{xurl}}{} % add URL line breaks if available
\urlstyle{same}
\hypersetup{
  pdftitle={Untitled},
  pdfauthor={kiki la petite sorcière},
  hidelinks,
  pdfcreator={LaTeX via pandoc}}

\title{Untitled}
\author{kiki la petite sorcière}
\date{2026-04-20}

\begin{document}
\maketitle

\section{Quelles cartes sont orientées attaque ou défense ? / Question
2}\label{quelles-cartes-sont-orientuxe9es-attaque-ou-duxe9fense-question-2}

Les cartes ne sont pas toujours équilibrées pour les attaquants ou les
défenseurs. Si l'une d'entre elles favorise les défenseurs, on pourra
expliquer en partie leurs victoires répétées. On étudiera d'abord pour
les joueurs normaux, puis on verra si les pros sont aussi sujets à ces
déséquilibres.

\subsection{Tous les joueurs}\label{tous-les-joueurs}

Commençons par étudier les données concernant les joueurs lambda grâce à
\href{https://blitz.gg/valorant/stats/maps?act=ev25actvi}{cette
ressource} qui nous donne les taux de victoire pour les côtés attaquant
et défenseur sur chaque map pour toute une saison de 2025.

Faisons une visualisation:

\begin{Shaded}
\begin{Highlighting}[]
\FunctionTok{ggplot}\NormalTok{(donnees, }\FunctionTok{aes}\NormalTok{(}\AttributeTok{x =}\NormalTok{ Map, }\AttributeTok{y =}\NormalTok{ WinRate, }\AttributeTok{fill =}\NormalTok{ Side)) }\SpecialCharTok{+}
  \FunctionTok{geom\_col}\NormalTok{(}\AttributeTok{width =} \FloatTok{0.75}\NormalTok{, }\AttributeTok{position =} \StringTok{"dodge"}\NormalTok{) }\SpecialCharTok{+}
  \FunctionTok{scale\_y\_continuous}\NormalTok{(}\AttributeTok{breaks =} \FunctionTok{seq}\NormalTok{(}\DecValTok{0}\NormalTok{, }\DecValTok{100}\NormalTok{, }\AttributeTok{by =} \DecValTok{5}\NormalTok{)) }\SpecialCharTok{+}
  \FunctionTok{theme}\NormalTok{(}\AttributeTok{axis.text.x =} \FunctionTok{element\_text}\NormalTok{(}\AttributeTok{angle =} \DecValTok{30}\NormalTok{, }\AttributeTok{hjust =} \DecValTok{1}\NormalTok{, }\AttributeTok{vjust =}\NormalTok{ .}\DecValTok{5}\NormalTok{, }\AttributeTok{face =} \StringTok{"bold"}\NormalTok{))}
\end{Highlighting}
\end{Shaded}

\pandocbounded{\includegraphics[keepaspectratio]{rapport_files/figure-latex/visu2-1.pdf}}

Finalement, les maps disponibles semblent assez équilibrées car toutes
les valeurs sont entre 45 et 55\%, même si on peut tout de même
constater quelques écarts. On remarque que la map Bind a le plus gros
écart, avec un winrate en défense proche de 46\%. A l'inverse, la map
Abyss semble la plus équilibrée ici, suivie de près par Corrode. Elles
semblent avoir tous leurs winrates entre 49 et 51\%.

\subsection{Joueurs professionnels}\label{joueurs-professionnels}

Voyons maintenant si les joueurs professsionnels ont des résultats
similaires, ou s'ils ont peut-être plus d'écarts.

Ici on utilise un tableau relatif aux stats des maps de l'année 2025 des
compétitions ESports.

Faisons une visualisation similaire à la première:

\begin{Shaded}
\begin{Highlighting}[]
\FunctionTok{ggplot}\NormalTok{(data, }\FunctionTok{aes}\NormalTok{(}\AttributeTok{x =}\NormalTok{ Map, }\AttributeTok{y =}\NormalTok{ WinRate, }\AttributeTok{fill =}\NormalTok{ Side)) }\SpecialCharTok{+}
  \FunctionTok{geom\_col}\NormalTok{(}\AttributeTok{width =} \FloatTok{0.75}\NormalTok{, }\AttributeTok{position =} \StringTok{"dodge"}\NormalTok{) }\SpecialCharTok{+}
  \FunctionTok{scale\_y\_continuous}\NormalTok{(}\AttributeTok{breaks =} \FunctionTok{seq}\NormalTok{(}\DecValTok{0}\NormalTok{, }\DecValTok{100}\NormalTok{, }\AttributeTok{by =} \DecValTok{5}\NormalTok{)) }\SpecialCharTok{+}
  \FunctionTok{theme}\NormalTok{(}\AttributeTok{axis.text.x =} \FunctionTok{element\_text}\NormalTok{(}\AttributeTok{angle =} \DecValTok{30}\NormalTok{, }\AttributeTok{hjust =} \DecValTok{1}\NormalTok{, }\AttributeTok{vjust =}\NormalTok{ .}\DecValTok{5}\NormalTok{, }\AttributeTok{face =} \StringTok{"bold"}\NormalTok{))}
\end{Highlighting}
\end{Shaded}

\pandocbounded{\includegraphics[keepaspectratio]{rapport_files/figure-latex/visu3-1.pdf}}

Etonnamment, la map Abyss qui semblait être la plus équilibrée pour les
joueurs lambda, s'avère être bien plus largement déséquilibrée chez les
pros, avec des winrates au dehors de la fourchette 45-55\%. Globalement,
il semble y avoir beaucoup plus de déséquilibres sur la scène sportive
que dans le jeu en général. C'est un résultat surprenant, mais on peut
imaginer que c'est parce que les joueurs professionnels ont des
connaissances bien plus avancées et des tactiques très précises selon la
map sur laquelle ils jouent qui leur permettent de tirer profit d'une
position ou d'une autre au maximum. Par ailleurs, les stats concernant
toutes les parties comprennent plusieurs millions de parties de joueurs
de niveaux très différents, tandis que les données sur les joueurs pro
concernent infiniment moins de parties jouées et de participants. Les
résultats sont donc peut-être biaisés.

\section{Qui sont les joueurs les plus forts en clutch ? / Question
4}\label{qui-sont-les-joueurs-les-plus-forts-en-clutch-question-4}

Être capable de retourner une situation désespérée est un atout
considérable pour un joueur professionnel de Valorant. Le clutch est le
fait de remporter une manche en étant le dernier membre de son équipe,
tout en devant se battre contre plusieurs adversaires. Nous allons voir
quels joueurs se retrouvent le plus souvent en situation de clutch (seul
contre tous), et lesquels s'en sortent le mieux. Nous utiliseront la
\href{https://www.kaggle.com/datasets/ryanluong1/valorant-champion-tour-2021-2023-data}{source
n°3}

\begin{Shaded}
\begin{Highlighting}[]
\NormalTok{top5 }\OtherTok{\textless{}{-}}\NormalTok{ df }\SpecialCharTok{\%\textgreater{}\%}
  \FunctionTok{group\_by}\NormalTok{(annee, Player) }\SpecialCharTok{\%\textgreater{}\%}
  \FunctionTok{summarise}\NormalTok{(}
    \StringTok{\textasciigrave{}}\AttributeTok{Clutches played}\StringTok{\textasciigrave{}} \OtherTok{=} \FunctionTok{sum}\NormalTok{(}\StringTok{\textasciigrave{}}\AttributeTok{Clutches played}\StringTok{\textasciigrave{}}\NormalTok{, }\AttributeTok{na.rm =} \ConstantTok{TRUE}\NormalTok{),}
    \StringTok{\textasciigrave{}}\AttributeTok{Clutches won}\StringTok{\textasciigrave{}} \OtherTok{=} \FunctionTok{sum}\NormalTok{(}\StringTok{\textasciigrave{}}\AttributeTok{Clutches won}\StringTok{\textasciigrave{}}\NormalTok{, }\AttributeTok{na.rm =} \ConstantTok{TRUE}\NormalTok{),}
    \AttributeTok{.groups =} \StringTok{"drop"}
\NormalTok{  ) }\SpecialCharTok{\%\textgreater{}\%}
  \FunctionTok{group\_by}\NormalTok{(annee) }\SpecialCharTok{\%\textgreater{}\%}
  \FunctionTok{slice\_max}\NormalTok{(}\StringTok{\textasciigrave{}}\AttributeTok{Clutches played}\StringTok{\textasciigrave{}}\NormalTok{, }\AttributeTok{n =} \DecValTok{5}\NormalTok{, }\AttributeTok{with\_ties =} \ConstantTok{FALSE}\NormalTok{) }\SpecialCharTok{\%\textgreater{}\%}
  \FunctionTok{ungroup}\NormalTok{()}

\NormalTok{plot\_year }\OtherTok{\textless{}{-}} \ControlFlowTok{function}\NormalTok{(y) \{}
\NormalTok{  df }\OtherTok{\textless{}{-}}\NormalTok{ top5 }\SpecialCharTok{\%\textgreater{}\%} \FunctionTok{filter}\NormalTok{(annee }\SpecialCharTok{==}\NormalTok{ y)}
  \FunctionTok{ggplot}\NormalTok{(df, }\FunctionTok{aes}\NormalTok{(}\AttributeTok{x =} \FunctionTok{reorder}\NormalTok{(Player, }\SpecialCharTok{{-}}\StringTok{\textasciigrave{}}\AttributeTok{Clutches played}\StringTok{\textasciigrave{}}\NormalTok{))) }\SpecialCharTok{+}
    \FunctionTok{geom\_col}\NormalTok{(}\FunctionTok{aes}\NormalTok{(}\AttributeTok{y =} \StringTok{\textasciigrave{}}\AttributeTok{Clutches played}\StringTok{\textasciigrave{}}\NormalTok{, }\AttributeTok{fill =} \StringTok{"Clutches perdus"}\NormalTok{), }\AttributeTok{fill =} \StringTok{"steelblue"}\NormalTok{) }\SpecialCharTok{+}
    \FunctionTok{geom\_col}\NormalTok{(}\FunctionTok{aes}\NormalTok{(}\AttributeTok{y =} \StringTok{\textasciigrave{}}\AttributeTok{Clutches won}\StringTok{\textasciigrave{}}\NormalTok{, }\AttributeTok{fill =} \StringTok{"Clutches gagnés"}\NormalTok{), }\AttributeTok{fill =} \StringTok{"gold"}\NormalTok{)}\SpecialCharTok{+}
    \FunctionTok{labs}\NormalTok{(}
      \AttributeTok{title =} \FunctionTok{paste}\NormalTok{(}\StringTok{"Top 5 joueurs qui clutchent le plus en"}\NormalTok{, y),}
      \AttributeTok{x =} \StringTok{"Joueur"}\NormalTok{,}
      \AttributeTok{y =} \StringTok{"Nombre de clutchs"}
\NormalTok{    ) }\SpecialCharTok{+}
    \FunctionTok{theme\_minimal}\NormalTok{()}
\NormalTok{\}}

\FunctionTok{plot\_year}\NormalTok{(}\DecValTok{2021}\NormalTok{)}
\end{Highlighting}
\end{Shaded}

\pandocbounded{\includegraphics[keepaspectratio]{rapport_files/figure-latex/visu1-1.pdf}}

\begin{Shaded}
\begin{Highlighting}[]
\FunctionTok{plot\_year}\NormalTok{(}\DecValTok{2022}\NormalTok{)}
\end{Highlighting}
\end{Shaded}

\pandocbounded{\includegraphics[keepaspectratio]{rapport_files/figure-latex/visu1-2.pdf}}

\begin{Shaded}
\begin{Highlighting}[]
\FunctionTok{plot\_year}\NormalTok{(}\DecValTok{2023}\NormalTok{)}
\end{Highlighting}
\end{Shaded}

\pandocbounded{\includegraphics[keepaspectratio]{rapport_files/figure-latex/visu1-3.pdf}}

\begin{Shaded}
\begin{Highlighting}[]
\FunctionTok{plot\_year}\NormalTok{(}\DecValTok{2024}\NormalTok{)}
\end{Highlighting}
\end{Shaded}

\pandocbounded{\includegraphics[keepaspectratio]{rapport_files/figure-latex/visu1-4.pdf}}

\begin{Shaded}
\begin{Highlighting}[]
\FunctionTok{plot\_year}\NormalTok{(}\DecValTok{2025}\NormalTok{)}
\end{Highlighting}
\end{Shaded}

\pandocbounded{\includegraphics[keepaspectratio]{rapport_files/figure-latex/visu1-5.pdf}}

\begin{Shaded}
\begin{Highlighting}[]
\FunctionTok{plot\_year}\NormalTok{(}\DecValTok{2026}\NormalTok{)}
\end{Highlighting}
\end{Shaded}

\pandocbounded{\includegraphics[keepaspectratio]{rapport_files/figure-latex/visu1-6.pdf}}

\subsection{Interprétation}\label{interpruxe9tation}

On a beaucoup d'informations intéressantes sur ces graphes. Déjà, on
voit que d'une année à l'autre, les joueurs ne sont pas les mêmes. On
voit donc qu'il n'existe pas de joueur spécialement meilleur que les
autres dans ces situations. On voit également que la quantité de clutchs
varie grandement d'une année à l'autre. Ce n'est pas parce que ces
situations sont plus ou moins rares, mais parce que dans certaines
années, il y a eu moins de matchs. En 2021, car le jeu venait de sortir,
et donc la scène compétitive comptait moins de joueurs, et en 2025 et
2026 car le dataset ne répertorie pas tous les matchs de ces années. On
remarque cependant un point commun chez tous les joueurs présents dans
ce tableau : ils jouent principalement des contrôleurs ou sentinelles,
c'est à dire des agents servant de soutien aux duellistes. On peut
expliquer cela car ces personnages sont souvent placés en retrait, et
survivent généralement plus longtemps que ceux qui cherchent à se
battre.

\section{Quels joueurs pros ont la carrière la plus prolifique
?}\label{quels-joueurs-pros-ont-la-carriuxe8re-la-plus-prolifique}

\subsection{Introduction}\label{introduction}

Dans cette partie, nous allons tenter de répondre à la question suivante
: \textbf{Quels joueurs pros ont la carrière la plus prolifique ?}
L'idée de cette question est d'observer les joueurs prometteurs pour
voir où est ce qu'on les retrouvera dans d'autres analyses.

\subsection{Démarche envisagée}\label{duxe9marche-envisaguxe9e}

Dans l'idée, nous souhaitons réaliser un barchart permettant de voir les
10 meilleurs joueurs de chaque année, afin d'en observer l'évolution.

\subsection{Réalisation}\label{ruxe9alisation}

\subsubsection{Chargement des
bibliothèques}\label{chargement-des-bibliothuxe8ques}

\begin{Shaded}
\begin{Highlighting}[]
\FunctionTok{library}\NormalTok{(ggplot2)}
\FunctionTok{library}\NormalTok{(magrittr)}
\FunctionTok{library}\NormalTok{ (dplyr)}
\FunctionTok{library}\NormalTok{ (readr)}
\end{Highlighting}
\end{Shaded}

\subsubsection{Chargement des datasets}\label{chargement-des-datasets}

\begin{Shaded}
\begin{Highlighting}[]
\NormalTok{perfs\_2021 }\OtherTok{\textless{}{-}} \FunctionTok{read\_csv}\NormalTok{(}\StringTok{"data/3. Valorant Champion Tour 2021{-}2026 Data/vct\_2021/players\_stats/players\_stats.csv"}\NormalTok{)}
\end{Highlighting}
\end{Shaded}

\begin{verbatim}
## Rows: 194966 Columns: 25
## -- Column specification --------------------------------------------------------
## Delimiter: ","
## chr (10): Tournament, Stage, Match Type, Player, Teams, Agents, Kill, Assist...
## dbl (15): Rounds Played, Rating, Average Combat Score, Kills:Deaths, Average...
## 
## i Use `spec()` to retrieve the full column specification for this data.
## i Specify the column types or set `show_col_types = FALSE` to quiet this message.
\end{verbatim}

\begin{Shaded}
\begin{Highlighting}[]
\NormalTok{perfs\_2022 }\OtherTok{\textless{}{-}} \FunctionTok{read\_csv}\NormalTok{(}\StringTok{"data/3. Valorant Champion Tour 2021{-}2026 Data/vct\_2022/players\_stats/players\_stats.csv"}\NormalTok{)}
\end{Highlighting}
\end{Shaded}

\begin{verbatim}
## Rows: 126641 Columns: 25
## -- Column specification --------------------------------------------------------
## Delimiter: ","
## chr (10): Tournament, Stage, Match Type, Player, Teams, Agents, Kill, Assist...
## dbl (15): Rounds Played, Rating, Average Combat Score, Kills:Deaths, Average...
## 
## i Use `spec()` to retrieve the full column specification for this data.
## i Specify the column types or set `show_col_types = FALSE` to quiet this message.
\end{verbatim}

\begin{Shaded}
\begin{Highlighting}[]
\NormalTok{perfs\_2023 }\OtherTok{\textless{}{-}} \FunctionTok{read\_csv}\NormalTok{(}\StringTok{"data/3. Valorant Champion Tour 2021{-}2026 Data/vct\_2023/players\_stats/players\_stats.csv"}\NormalTok{)}
\end{Highlighting}
\end{Shaded}

\begin{verbatim}
## Rows: 11249 Columns: 25
## -- Column specification --------------------------------------------------------
## Delimiter: ","
## chr (10): Tournament, Stage, Match Type, Player, Teams, Agents, Kill, Assist...
## dbl (15): Rounds Played, Rating, Average Combat Score, Kills:Deaths, Average...
## 
## i Use `spec()` to retrieve the full column specification for this data.
## i Specify the column types or set `show_col_types = FALSE` to quiet this message.
\end{verbatim}

\begin{Shaded}
\begin{Highlighting}[]
\NormalTok{perfs\_2024 }\OtherTok{\textless{}{-}} \FunctionTok{read\_csv}\NormalTok{(}\StringTok{"data/3. Valorant Champion Tour 2021{-}2026 Data/vct\_2024/players\_stats/players\_stats.csv"}\NormalTok{)}
\end{Highlighting}
\end{Shaded}

\begin{verbatim}
## Rows: 15030 Columns: 25
## -- Column specification --------------------------------------------------------
## Delimiter: ","
## chr (10): Tournament, Stage, Match Type, Player, Teams, Agents, Kill, Assist...
## dbl (15): Rounds Played, Rating, Average Combat Score, Kills:Deaths, Average...
## 
## i Use `spec()` to retrieve the full column specification for this data.
## i Specify the column types or set `show_col_types = FALSE` to quiet this message.
\end{verbatim}

\begin{Shaded}
\begin{Highlighting}[]
\NormalTok{perfs\_2025 }\OtherTok{\textless{}{-}} \FunctionTok{read\_csv}\NormalTok{(}\StringTok{"data/3. Valorant Champion Tour 2021{-}2026 Data/vct\_2025/players\_stats/players\_stats.csv"}\NormalTok{)}
\end{Highlighting}
\end{Shaded}

\begin{verbatim}
## Rows: 17996 Columns: 25
## -- Column specification --------------------------------------------------------
## Delimiter: ","
## chr (10): Tournament, Stage, Match Type, Player, Teams, Agents, Kill, Assist...
## dbl (15): Rounds Played, Rating, Average Combat Score, Kills:Deaths, Average...
## 
## i Use `spec()` to retrieve the full column specification for this data.
## i Specify the column types or set `show_col_types = FALSE` to quiet this message.
\end{verbatim}

\subsubsection{Réalisation des
graphiques}\label{ruxe9alisation-des-graphiques}

\begin{Shaded}
\begin{Highlighting}[]
\CommentTok{\# Graphique 2021}
\NormalTok{unique\_players\_2021 }\OtherTok{\textless{}{-}}\NormalTok{ perfs\_2021 }\SpecialCharTok{\%\textgreater{}\%} \FunctionTok{filter}\NormalTok{(}\SpecialCharTok{!}\FunctionTok{duplicated}\NormalTok{(Player)) }\CommentTok{\#Suppression des doublons}
\NormalTok{top\_players\_2021 }\OtherTok{\textless{}{-}} \FunctionTok{head}\NormalTok{(unique\_players\_2021 }\SpecialCharTok{\%\textgreater{}\%} \FunctionTok{arrange}\NormalTok{(}\FunctionTok{desc}\NormalTok{(Rating)), }\DecValTok{10}\NormalTok{) }\CommentTok{\# Récupération des meilleurs joueurs}
\FunctionTok{ggplot}\NormalTok{(top\_players\_2021, }\FunctionTok{aes}\NormalTok{(}\AttributeTok{x =} \FunctionTok{reorder}\NormalTok{(Player, }\SpecialCharTok{{-}}\NormalTok{Rating), }\AttributeTok{y =}\NormalTok{ Rating)) }\SpecialCharTok{+} \FunctionTok{geom\_bar}\NormalTok{(}\AttributeTok{stat =} \StringTok{"identity"}\NormalTok{) }\SpecialCharTok{+} \FunctionTok{labs}\NormalTok{(}\AttributeTok{title =} \StringTok{"Bar chart des meilleurs joueurs de 2021"}\NormalTok{, }\AttributeTok{x =} \StringTok{"Joueur"}\NormalTok{, }\AttributeTok{y =} \StringTok{"Rating"}\NormalTok{) }\CommentTok{\#Graphique}
\end{Highlighting}
\end{Shaded}

\pandocbounded{\includegraphics[keepaspectratio]{rapport_files/figure-latex/unnamed-chunk-3-1.pdf}}

\begin{Shaded}
\begin{Highlighting}[]
\CommentTok{\# Graphique 2022}
\NormalTok{unique\_players\_2022 }\OtherTok{\textless{}{-}}\NormalTok{ perfs\_2022 }\SpecialCharTok{\%\textgreater{}\%} \FunctionTok{filter}\NormalTok{(}\SpecialCharTok{!}\FunctionTok{duplicated}\NormalTok{(Player)) }\CommentTok{\#Suppression des doublons}
\NormalTok{top\_players\_2022 }\OtherTok{\textless{}{-}} \FunctionTok{head}\NormalTok{(unique\_players\_2022 }\SpecialCharTok{\%\textgreater{}\%} \FunctionTok{arrange}\NormalTok{(}\FunctionTok{desc}\NormalTok{(Rating)), }\DecValTok{10}\NormalTok{) }\CommentTok{\# Récupération des meilleurs joueurs}
\FunctionTok{ggplot}\NormalTok{(top\_players\_2022, }\FunctionTok{aes}\NormalTok{(}\AttributeTok{x =} \FunctionTok{reorder}\NormalTok{(Player, }\SpecialCharTok{{-}}\NormalTok{Rating), }\AttributeTok{y =}\NormalTok{ Rating)) }\SpecialCharTok{+} \FunctionTok{geom\_bar}\NormalTok{(}\AttributeTok{stat =} \StringTok{"identity"}\NormalTok{) }\SpecialCharTok{+} \FunctionTok{labs}\NormalTok{(}\AttributeTok{title =} \StringTok{"Bar chart des meilleurs joueurs de 2022"}\NormalTok{, }\AttributeTok{x =} \StringTok{"Joueur"}\NormalTok{, }\AttributeTok{y =} \StringTok{"Rating"}\NormalTok{) }\CommentTok{\#Graphique}
\end{Highlighting}
\end{Shaded}

\pandocbounded{\includegraphics[keepaspectratio]{rapport_files/figure-latex/unnamed-chunk-4-1.pdf}}

\begin{Shaded}
\begin{Highlighting}[]
\CommentTok{\# Graphique 2023}
\NormalTok{unique\_players\_2023 }\OtherTok{\textless{}{-}}\NormalTok{ perfs\_2023 }\SpecialCharTok{\%\textgreater{}\%} \FunctionTok{filter}\NormalTok{(}\SpecialCharTok{!}\FunctionTok{duplicated}\NormalTok{(Player)) }\CommentTok{\#Suppression des doublons}
\NormalTok{top\_players\_2023 }\OtherTok{\textless{}{-}} \FunctionTok{head}\NormalTok{(unique\_players\_2023 }\SpecialCharTok{\%\textgreater{}\%} \FunctionTok{arrange}\NormalTok{(}\FunctionTok{desc}\NormalTok{(Rating)), }\DecValTok{10}\NormalTok{) }\CommentTok{\# Récupération des meilleurs joueurs}
\FunctionTok{ggplot}\NormalTok{(top\_players\_2023, }\FunctionTok{aes}\NormalTok{(}\AttributeTok{x =} \FunctionTok{reorder}\NormalTok{(Player, }\SpecialCharTok{{-}}\NormalTok{Rating), }\AttributeTok{y =}\NormalTok{ Rating)) }\SpecialCharTok{+} \FunctionTok{geom\_bar}\NormalTok{(}\AttributeTok{stat =} \StringTok{"identity"}\NormalTok{) }\SpecialCharTok{+} \FunctionTok{labs}\NormalTok{(}\AttributeTok{title =} \StringTok{"Bar chart des meilleurs joueurs de 2023"}\NormalTok{, }\AttributeTok{x =} \StringTok{"Joueur"}\NormalTok{, }\AttributeTok{y =} \StringTok{"Rating"}\NormalTok{)}\CommentTok{\#Graphique}
\end{Highlighting}
\end{Shaded}

\pandocbounded{\includegraphics[keepaspectratio]{rapport_files/figure-latex/unnamed-chunk-5-1.pdf}}

\begin{Shaded}
\begin{Highlighting}[]
\CommentTok{\# Graphique 2024}
\NormalTok{unique\_players\_2024 }\OtherTok{\textless{}{-}}\NormalTok{ perfs\_2024 }\SpecialCharTok{\%\textgreater{}\%} \FunctionTok{filter}\NormalTok{(}\SpecialCharTok{!}\FunctionTok{duplicated}\NormalTok{(Player)) }\CommentTok{\#Suppression des doublons}
\NormalTok{top\_players\_2024 }\OtherTok{\textless{}{-}} \FunctionTok{head}\NormalTok{(unique\_players\_2024 }\SpecialCharTok{\%\textgreater{}\%} \FunctionTok{arrange}\NormalTok{(}\FunctionTok{desc}\NormalTok{(Rating)), }\DecValTok{10}\NormalTok{) }\CommentTok{\# Récupération des meilleurs joueurs}
\FunctionTok{ggplot}\NormalTok{(top\_players\_2024, }\FunctionTok{aes}\NormalTok{(}\AttributeTok{x =} \FunctionTok{reorder}\NormalTok{(Player, }\SpecialCharTok{{-}}\NormalTok{Rating), }\AttributeTok{y =}\NormalTok{ Rating)) }\SpecialCharTok{+} \FunctionTok{geom\_bar}\NormalTok{(}\AttributeTok{stat =} \StringTok{"identity"}\NormalTok{) }\SpecialCharTok{+} \FunctionTok{labs}\NormalTok{(}\AttributeTok{title =} \StringTok{"Bar chart des meilleurs joueurs de 2024"}\NormalTok{, }\AttributeTok{x =} \StringTok{"Joueur"}\NormalTok{, }\AttributeTok{y =} \StringTok{"Rating"}\NormalTok{) }\CommentTok{\#Graphique}
\end{Highlighting}
\end{Shaded}

\pandocbounded{\includegraphics[keepaspectratio]{rapport_files/figure-latex/unnamed-chunk-6-1.pdf}}

\begin{Shaded}
\begin{Highlighting}[]
\CommentTok{\# Graphique 2025}
\NormalTok{unique\_players\_2025 }\OtherTok{\textless{}{-}}\NormalTok{ perfs\_2025 }\SpecialCharTok{\%\textgreater{}\%} \FunctionTok{filter}\NormalTok{(}\SpecialCharTok{!}\FunctionTok{duplicated}\NormalTok{(Player)) }\CommentTok{\#Suppression des doublons}
\NormalTok{top\_players\_2025 }\OtherTok{\textless{}{-}} \FunctionTok{head}\NormalTok{(unique\_players\_2025 }\SpecialCharTok{\%\textgreater{}\%} \FunctionTok{arrange}\NormalTok{(}\FunctionTok{desc}\NormalTok{(Rating)), }\DecValTok{10}\NormalTok{) }\CommentTok{\# Récupération des meilleurs joueurs}
\FunctionTok{ggplot}\NormalTok{(top\_players\_2025, }\FunctionTok{aes}\NormalTok{(}\AttributeTok{x =} \FunctionTok{reorder}\NormalTok{(Player, }\SpecialCharTok{{-}}\NormalTok{Rating), }\AttributeTok{y =}\NormalTok{ Rating)) }\SpecialCharTok{+} \FunctionTok{geom\_bar}\NormalTok{(}\AttributeTok{stat =} \StringTok{"identity"}\NormalTok{) }\SpecialCharTok{+}
\FunctionTok{labs}\NormalTok{(}\AttributeTok{title =} \StringTok{"Bar chart des meilleurs joueurs de 2025"}\NormalTok{, }\AttributeTok{x =} \StringTok{"Joueur"}\NormalTok{, }\AttributeTok{y =} \StringTok{"Rating"}\NormalTok{) }\CommentTok{\#Graphique}
\end{Highlighting}
\end{Shaded}

\pandocbounded{\includegraphics[keepaspectratio]{rapport_files/figure-latex/unnamed-chunk-7-1.pdf}}

\subsubsection{Line Chart de l'évolution des meilleurs
joueurs}\label{line-chart-de-luxe9volution-des-meilleurs-joueurs}

Voyant que les 10 meilleurs joueurs changent régulièrement, nous avons
souhaité vérifier si les meilleurs joueurs restent dans le top 10
plusieurs années d'affilée

\begin{Shaded}
\begin{Highlighting}[]
\NormalTok{top\_players }\OtherTok{\textless{}{-}} \FunctionTok{bind\_rows}\NormalTok{(}
\NormalTok{  top\_players\_2021 }\SpecialCharTok{\%\textgreater{}\%} \FunctionTok{mutate}\NormalTok{(}\AttributeTok{annee =} \DecValTok{2021}\NormalTok{),}
\NormalTok{  top\_players\_2022 }\SpecialCharTok{\%\textgreater{}\%} \FunctionTok{mutate}\NormalTok{(}\AttributeTok{annee =} \DecValTok{2022}\NormalTok{),}
\NormalTok{  top\_players\_2023 }\SpecialCharTok{\%\textgreater{}\%} \FunctionTok{mutate}\NormalTok{(}\AttributeTok{annee =} \DecValTok{2023}\NormalTok{),}
\NormalTok{  top\_players\_2024 }\SpecialCharTok{\%\textgreater{}\%} \FunctionTok{mutate}\NormalTok{(}\AttributeTok{annee =} \DecValTok{2024}\NormalTok{),}
\NormalTok{  top\_players\_2025 }\SpecialCharTok{\%\textgreater{}\%} \FunctionTok{mutate}\NormalTok{(}\AttributeTok{annee =} \DecValTok{2025}\NormalTok{)}
\NormalTok{)}
\end{Highlighting}
\end{Shaded}

\begin{Shaded}
\begin{Highlighting}[]
\FunctionTok{ggplot}\NormalTok{(top\_players, }\FunctionTok{aes}\NormalTok{(}\AttributeTok{x =}\NormalTok{ annee, }\AttributeTok{y =}\NormalTok{ Rating, }\AttributeTok{group =}\NormalTok{ Player)) }\SpecialCharTok{+} \FunctionTok{geom\_line}\NormalTok{(}\FunctionTok{aes}\NormalTok{(}\AttributeTok{color=}\NormalTok{Player)) }\SpecialCharTok{+} \FunctionTok{geom\_point}\NormalTok{(}\FunctionTok{aes}\NormalTok{(}\AttributeTok{color=}\NormalTok{Player)) }\SpecialCharTok{+}
\FunctionTok{labs}\NormalTok{(}\AttributeTok{title =} \StringTok{"Bar chart des meilleurs joueurs de 2025"}\NormalTok{, }\AttributeTok{x =} \StringTok{"Joueur"}\NormalTok{, }\AttributeTok{y =} \StringTok{"Rating"}\NormalTok{, }\AttributeTok{colour =} \StringTok{"Player"}\NormalTok{) }\CommentTok{\#Graphique}
\end{Highlighting}
\end{Shaded}

\begin{verbatim}
## `geom_line()`: Each group consists of only one observation.
## i Do you need to adjust the group aesthetic?
\end{verbatim}

\pandocbounded{\includegraphics[keepaspectratio]{rapport_files/figure-latex/unnamed-chunk-9-1.pdf}}

\subsection{Observations}\label{observations}

Le choix de représentation nous amène à observer que les premiers
joueurs changent chaque année et que donc il est très difficile
d'observer la progression de seulement les meilleurs joueurs. Le dernier
graphique montre d'ailleurs que chaque joueur n'apparait jamais
plusieurs fois dans le top 10. Il est donc impossible d'en percevoir
l'évolution.

\section{Quelles sont les meilleures équipes par année
?}\label{quelles-sont-les-meilleures-uxe9quipes-par-annuxe9e}

Complémentaire à la question précédente, on pourra ainsi voir
l'évolution des équipes, et la corréler aux changements de joueurs ou de
coachs, ou bien d'ajouts de cartes ou d'agents, pour voir les causes de
ces évolutions.

\subsection{Import des données}\label{import-des-donnuxe9es}

On importe les données des Valorant Champion 2021-2025 en récupérant:

\begin{itemize}
\item
  Le winrate par équipe
\item
  Le nombre de changements au roster entre 2022 et 2025
\end{itemize}

\begin{Shaded}
\begin{Highlighting}[]
\NormalTok{resultat\_final }\OtherTok{\textless{}{-}} \FunctionTok{map\_df}\NormalTok{(}\DecValTok{2021}\SpecialCharTok{:}\DecValTok{2025}\NormalTok{, \textbackslash{}(an) \{}
  \FunctionTok{read.csv}\NormalTok{(}\FunctionTok{sprintf}\NormalTok{(}\StringTok{"data/3. Valorant Champion Tour 2021{-}2026 Data/vct\_\%d/matches/overview.csv"}\NormalTok{, an)) }\SpecialCharTok{\%\textgreater{}\%}
    \FunctionTok{filter}\NormalTok{(Tournament }\SpecialCharTok{==} \FunctionTok{paste}\NormalTok{(}\StringTok{"Valorant Champions"}\NormalTok{, an)) }\SpecialCharTok{\%\textgreater{}\%}
    \FunctionTok{distinct}\NormalTok{(}\AttributeTok{Annee =}\NormalTok{ an, }\AttributeTok{Equipe =}\NormalTok{ Team, Player)}
\NormalTok{\})}

\NormalTok{roster\_data }\OtherTok{\textless{}{-}}\NormalTok{ resultat\_final}

\NormalTok{resultat\_final }\OtherTok{\textless{}{-}}\NormalTok{ roster\_data }\SpecialCharTok{\%\textgreater{}\%}
  \FunctionTok{group\_by}\NormalTok{(Equipe) }\SpecialCharTok{\%\textgreater{}\%}
  \FunctionTok{arrange}\NormalTok{(Annee) }\SpecialCharTok{\%\textgreater{}\%}
  \FunctionTok{mutate}\NormalTok{(}\AttributeTok{Changements =} \FunctionTok{map\_int}\NormalTok{(Annee, \textbackslash{}(a) \{}
    \ControlFlowTok{if}\NormalTok{ (a }\SpecialCharTok{==} \DecValTok{2021}\NormalTok{) }\FunctionTok{return}\NormalTok{(}\DecValTok{0}\NormalTok{)}
\NormalTok{    curr }\OtherTok{\textless{}{-}}\NormalTok{ roster\_data}\SpecialCharTok{$}\NormalTok{Player[roster\_data}\SpecialCharTok{$}\NormalTok{Annee }\SpecialCharTok{==}\NormalTok{ a }\SpecialCharTok{\&}\NormalTok{ roster\_data}\SpecialCharTok{$}\NormalTok{Equipe }\SpecialCharTok{==} \FunctionTok{first}\NormalTok{(Equipe)]}
\NormalTok{    prev }\OtherTok{\textless{}{-}}\NormalTok{ roster\_data}\SpecialCharTok{$}\NormalTok{Player[roster\_data}\SpecialCharTok{$}\NormalTok{Annee }\SpecialCharTok{==}\NormalTok{ a }\SpecialCharTok{{-}} \DecValTok{1} \SpecialCharTok{\&}\NormalTok{ roster\_data}\SpecialCharTok{$}\NormalTok{Equipe }\SpecialCharTok{==} \FunctionTok{first}\NormalTok{(Equipe)]}
    \ControlFlowTok{if}\NormalTok{ (}\FunctionTok{length}\NormalTok{(prev) }\SpecialCharTok{==} \DecValTok{0}\NormalTok{) }\FunctionTok{return}\NormalTok{(}\DecValTok{0}\NormalTok{)}
    \FunctionTok{length}\NormalTok{(}\FunctionTok{setdiff}\NormalTok{(prev, curr))}
\NormalTok{  \})) }\SpecialCharTok{\%\textgreater{}\%}
  \FunctionTok{distinct}\NormalTok{(Annee, Equipe, Changements) }\SpecialCharTok{\%\textgreater{}\%}
  \FunctionTok{ungroup}\NormalTok{()}

\NormalTok{tableau\_final }\OtherTok{\textless{}{-}} \FunctionTok{map\_df}\NormalTok{(}\DecValTok{2021}\SpecialCharTok{:}\DecValTok{2025}\NormalTok{, \textbackslash{}(an) \{}
\NormalTok{  path }\OtherTok{\textless{}{-}} \FunctionTok{sprintf}\NormalTok{(}\StringTok{"data/3. Valorant Champion Tour 2021{-}2026 Data/vct\_\%d/matches/scores.csv"}\NormalTok{, an)}
  \FunctionTok{read.csv}\NormalTok{(path) }\SpecialCharTok{\%\textgreater{}\%}
    \FunctionTok{filter}\NormalTok{(Tournament }\SpecialCharTok{==} \FunctionTok{paste}\NormalTok{(}\StringTok{"Valorant Champions"}\NormalTok{, an)) }\SpecialCharTok{\%\textgreater{}\%}
    \FunctionTok{reframe}\NormalTok{(}\AttributeTok{Team =} \FunctionTok{c}\NormalTok{(Team.A, Team.B), }\AttributeTok{S\_Self =} \FunctionTok{c}\NormalTok{(Team.A.Score, Team.B.Score), }\AttributeTok{S\_Opp =} \FunctionTok{c}\NormalTok{(Team.B.Score, Team.A.Score)) }\SpecialCharTok{\%\textgreater{}\%}
    \FunctionTok{group\_by}\NormalTok{(Team) }\SpecialCharTok{\%\textgreater{}\%}
    \FunctionTok{summarise}\NormalTok{(}\AttributeTok{WinRate =} \FunctionTok{mean}\NormalTok{(S\_Self }\SpecialCharTok{\textgreater{}}\NormalTok{ S\_Opp), Année }\OtherTok{=}\NormalTok{ an, }\AttributeTok{.groups =} \StringTok{"drop"}\NormalTok{)}
\NormalTok{\})}

\NormalTok{data\_combinee }\OtherTok{\textless{}{-}}\NormalTok{ tableau\_final }\SpecialCharTok{\%\textgreater{}\%}
  \FunctionTok{rename}\NormalTok{(}\AttributeTok{Equipe =}\NormalTok{ Team) }\SpecialCharTok{\%\textgreater{}\%}
  \FunctionTok{left\_join}\NormalTok{(resultat\_final, }\AttributeTok{by =} \FunctionTok{c}\NormalTok{(}\StringTok{"Année"} \OtherTok{=} \StringTok{"Annee"}\NormalTok{, }\StringTok{"Equipe"}\NormalTok{)) }\SpecialCharTok{\%\textgreater{}\%}
  \FunctionTok{mutate}\NormalTok{(}\AttributeTok{Changements =} \FunctionTok{replace\_na}\NormalTok{(Changements, }\DecValTok{0}\NormalTok{)) }\SpecialCharTok{\%\textgreater{}\%}
  \FunctionTok{filter}\NormalTok{(Equipe }\SpecialCharTok{\%in\%} \FunctionTok{names}\NormalTok{(}\FunctionTok{which}\NormalTok{(}\FunctionTok{table}\NormalTok{(Equipe) }\SpecialCharTok{\textgreater{}=} \DecValTok{3}\NormalTok{)))}
\end{Highlighting}
\end{Shaded}

Affichage des données

\begin{Shaded}
\begin{Highlighting}[]
\FunctionTok{ggplot}\NormalTok{(data\_combinee, }\FunctionTok{aes}\NormalTok{(}\AttributeTok{x =}\NormalTok{ Année, }\AttributeTok{group =}\NormalTok{ Equipe)) }\SpecialCharTok{+}
  \FunctionTok{geom\_area}\NormalTok{(}\FunctionTok{aes}\NormalTok{(}\AttributeTok{y =}\NormalTok{ Changements }\SpecialCharTok{/} \DecValTok{11}\NormalTok{), }\AttributeTok{fill =} \StringTok{"orchid4"}\NormalTok{, }\AttributeTok{alpha =} \FloatTok{0.2}\NormalTok{) }\SpecialCharTok{+} 
  \FunctionTok{geom\_line}\NormalTok{(}\FunctionTok{aes}\NormalTok{(}\AttributeTok{y =}\NormalTok{ Changements }\SpecialCharTok{/} \DecValTok{11}\NormalTok{), }\AttributeTok{color =} \StringTok{"orchid3"}\NormalTok{, }\AttributeTok{alpha =} \FloatTok{0.8}\NormalTok{) }\SpecialCharTok{+}
  \FunctionTok{geom\_line}\NormalTok{(}\FunctionTok{aes}\NormalTok{(}\AttributeTok{y =}\NormalTok{ WinRate, }\AttributeTok{color =}\NormalTok{ Equipe), }\AttributeTok{linewidth =} \DecValTok{1}\NormalTok{) }\SpecialCharTok{+}
  \FunctionTok{geom\_point}\NormalTok{(}\FunctionTok{aes}\NormalTok{(}\AttributeTok{y =}\NormalTok{ WinRate, }\AttributeTok{color =}\NormalTok{ Equipe), }\AttributeTok{size =} \FloatTok{2.5}\NormalTok{) }\SpecialCharTok{+}
  \FunctionTok{facet\_wrap}\NormalTok{(}\SpecialCharTok{\textasciitilde{}}\NormalTok{ Equipe) }\SpecialCharTok{+}
  \FunctionTok{scale\_y\_continuous}\NormalTok{(}
    \AttributeTok{name =} \StringTok{"Win Rate (\%)"}\NormalTok{,}
    \AttributeTok{limits =} \DecValTok{0}\SpecialCharTok{:}\DecValTok{1}\NormalTok{, }\AttributeTok{labels =}\NormalTok{ scales}\SpecialCharTok{::}\NormalTok{percent,}
    \AttributeTok{sec.axis =} \FunctionTok{sec\_axis}\NormalTok{(}\SpecialCharTok{\textasciitilde{}}\NormalTok{ . }\SpecialCharTok{*} \DecValTok{11}\NormalTok{, }\AttributeTok{name =} \StringTok{"Nombre de changements de joueurs"}\NormalTok{)}
\NormalTok{  ) }\SpecialCharTok{+}
  \FunctionTok{scale\_x\_continuous}\NormalTok{(}\AttributeTok{breaks =} \DecValTok{2021}\SpecialCharTok{:}\DecValTok{2025}\NormalTok{) }\SpecialCharTok{+}
  \FunctionTok{theme\_minimal}\NormalTok{() }\SpecialCharTok{+}
  \FunctionTok{theme}\NormalTok{(}\AttributeTok{legend.position =} \StringTok{"none"}\NormalTok{, }
        \AttributeTok{axis.title.y.right =} \FunctionTok{element\_text}\NormalTok{(}\AttributeTok{color =} \StringTok{"orchid4"}\NormalTok{, }\AttributeTok{size =} \DecValTok{9}\NormalTok{),}
        \AttributeTok{panel.grid.minor =} \FunctionTok{element\_blank}\NormalTok{()) }\SpecialCharTok{+}
  \FunctionTok{labs}\NormalTok{(}\AttributeTok{title =} \StringTok{"Évolution du winrate aux Valorant Champion (2021{-}2025, 3 participation min.)"}\NormalTok{)}
\end{Highlighting}
\end{Shaded}

\pandocbounded{\includegraphics[keepaspectratio]{rapport_files/figure-latex/unnamed-chunk-11-1.pdf}}

\subsection{Interprétation}\label{interpruxe9tation-1}

On a filtré les équipes pour n'afficher que les équipes avec le plus de
longévité aux Valorant Champions (plus de 3 participations sur
2021-2025). On a choisi de représenter le Winrate (Matchs gagnés/Matchs
joués) pour afficher directement la performance.

Ce que l'on remarque, c'est que certaines équipes arrivent à se
qualifier fréquemment, alors que la majorité des équipes qualifiées a
1-2 occurrences mais cela ne les empêche pas de remporter les Valorant
Champion tout de même. Concernant les équipes qui se qualifient
régulièrement, DRX, FNATIC et Paper Rex semblent se distinguer par des
performances qualitatives assez stables à travers les années

Le nombre de changement de joueurs semble indiquer que les changements
de joueurs ont pour conséquence de faire baisser les performances de
l'équipe, tandis qu'un taux de changement bas est souvent corrélé avec
des performances stables sur le long terme.

\end{document}
